\documentclass[11pt]{article}
\usepackage{amsmath, amssymb, booktabs, graphicx, xcolor}
\usepackage[margin=1in]{geometry}
\usepackage{bbm}          % in the preamble



\begin{document}

\section*{Technical Note: Distance to Government and Support for Technocratic Rule}

\subsection*{1. Constructing \texttt{gov\_distance\_w\_01}}

The main independent variable is \texttt{gov\_distance\_w\_01}, which captures the party-weighted distance from individual $i$ from the current Government.

Let $i$ index respondents and $p \in G$ index parties in the governing coalition.
From Q27, respondents rate their \emph{closeness} to each parliamentary party on a $0$--$10$ scale
($12$ = ``don't know''). I first set all $12$ to missing. 

Let $C_{ip} \in [0,10]$ denote respondent $i$'s closeness to party $p$.
Let $s_p$ be party $p$'s seat count in the Eduskunta at the time of fieldwork
(\emph{KOK} $=48$, \emph{PS} $=46$, \emph{RKP} $=9$, \emph{KD} $=5$),
and define the coalition seat weights
\[
w_p \equiv \frac{s_p}{\sum_{q \in G} s_q}.
\]
Because some respondents skip one or more government parties when answering Q27, I \emph{renormalize} the weights over the parties that respondent $i$ actually rated:
\[
\tilde w_{ip} \equiv 
\frac{w_p \cdot \mathbbm{1}\{C_{ip}\text{ is observed}\}}
{\sum_{q \in G} w_q \cdot \mathbbm{1}\{C_{iq}\ \text{is observed}\}}.
\]
The respondent's seat-weighted coalition \emph{closeness} is then
\[
\text{GovCloseness}_i \equiv \sum_{p \in G} \tilde w_{ip}\, C_{ip}
\quad \in [0,10].
\]
I scaled it to $0$--$1$ and defined a complementary \emph{distance} index:
\[
\texttt{gov\_closeness\_w\_01}_i \equiv \frac{\text{GovCloseness}_i}{10},
\qquad
\texttt{gov\_distance\_w\_01}_i \equiv 1 - \texttt{gov\_closeness\_w\_01}_i.
\]
By construction, higher values of \texttt{gov\_distance\_w\_01} indicate that a respondent is
farther from the governing coalition in the Q27 closeness space.

\subsection*{2. Estimation model}

I modeled support for technocratic delegation using a linear specification,
treating the $1$--$5$ item as approximately continuous (I also estimated an ordered model as a robustness check):
\[
\begin{aligned}
\underbrace{\texttt{techno5}_i}_{\text{Technocracy}}
&= \beta_0 + \beta_1\,\texttt{gov\_distance\_w\_01}_i
  + \beta_2\,\underbrace{\texttt{Q8\_1}_i}_{\text{Dem satisfaction}}
  + \beta_3\,\underbrace{\texttt{Q9\_4}_i}_{\text{Trust in Finnish politicians}} \\
&\quad + \boldsymbol{\gamma}'\,\mathbf{X}_i
  + \delta_{\text{region}(i)} + \varepsilon_i .
\end{aligned}
\]

where $\mathbf{X}_i$ includes education, age and gender; $\delta_{\text{region}}$ are region fixed effects. The coefficient of interest is $\beta_1$.

\subsection*{3. Interpretation of the predictive plot}

Figure~\ref{fig:pred} plots the \emph{average} predicted value of \texttt{techno5}
over the empirical distribution of controls, as \texttt{gov\_distance\_w\_01} varies from $0$ (very close to the government) to $1$ (far).
The dark line is the point prediction; the shaded band is the 95\% confidence interval.
The estimated relationship is \emph{negative and monotonic}: moving from full proximity to maximal distance lowers
predicted technocratic support from about $3.5$ to about $2.8$ on the $1$--$5$ scale.

\begin{figure}[h!]
\centering
\includegraphics[width=.8\linewidth]{Rplot.pdf} % replace with your path/filename
\caption{Average predicted technocracy vs.\ government distance (95\% CIs).}
\label{fig:pred}
\end{figure}

\subsection*{4. Substantive takeaway and relation to our expectation}

Our expectation was that \emph{electoral losers}---those farther from the sitting government---would be
more likely to (a) defend checks and balances and (b) support delegating power away from elected politicians (i.e., more technocracy). The evidence in Figure~\ref{fig:pred} goes in the opposite direction for (b): in Finland,
respondents who are \emph{farther} from the right-wing Orpo coalition exhibit \emph{lower} support for technocratic delegation.

A plausible interpretation is that, in this context, distance from the government largely identifies
left-leaning respondents. For these voters, ``experts rather than politicians'' may be perceived as
\emph{depoliticizing distributional conflict} and historically associated with fiscal austerity and
market-forward governance. Thus, when the right governs, left-leaning losers appear
\emph{less} willing to transfer authority to non-elected experts and prefer to keep decisions under elected, accountable politicians.
This pattern can coexist with stronger support for \emph{checks and balances} (judicial review, oversight),
which constrain incumbents without removing decisions from electoral accountability. In short,
losers may favor \emph{constraints on the government}, but not \emph{depoliticization} of policy.


\medskip
\noindent
\textit{Implications.} The results suggest that support for technocracy is not a generic ``losers' consent'' mechanism.
It appears contingent on the ideological meaning of expertise in a given polity: when expertise is coded
as center-right (Finland, 2023--2025), distance from the government correlates with \emph{lower} technocratic support.

\end{document}